% ============================================================
\chapter{Case Studies and Projects}
\label{ch:case_studies}
% ============================================================

\begin{center}
\textit{``All models are wrong, but some are useful.''}\\[4pt]
--- George E.\ P.\ Box
\end{center}

\bigskip

This chapter applies the tools developed in earlier chapters to five
concrete case studies. Each study illustrates the full modeling cycle:
formulation, mathematical analysis, calibration with data, and validation.

% ------------------------------------------------------------
\section{Epidemiological Modeling: SIR and SEIR}
\label{sec:sir_seir}
% ------------------------------------------------------------

\begin{definition}[SIR Model]
\index{SIR model}
The \textbf{SIR} (Susceptible--Infected--Recovered) model of
Kermack--McKendrick\index{Kermack--McKendrick} describes disease spread
in a population of size $N$:
\[
  \frac{dS}{dt} = -\beta\,\frac{SI}{N}, \quad
  \frac{dI}{dt} = \beta\,\frac{SI}{N} - \gamma\,I, \quad
  \frac{dR}{dt} = \gamma\,I,
\]
where $\beta > 0$ is the transmission rate and $\gamma > 0$ the recovery rate.
\end{definition}

\begin{theorem}[Epidemic Threshold]
\label{thm:threshold_sir}
Let $R_0 = \beta/\gamma$\index{basic reproduction number} be the
\emph{basic reproduction number}. If $R_0 > 1$ and $S(0) \approx N$,
then $I(t)$ initially increases (epidemic outbreak). If $R_0 < 1$,
the number of infected individuals decreases monotonically.
\end{theorem}

\begin{proof}
At the onset of the epidemic, $S \approx N$, so
$dI/dt \approx (\beta - \gamma)\,I = \gamma(R_0 - 1)\,I$.
If $R_0 > 1$, this linearized equation yields exponential growth of $I$.
If $R_0 < 1$, $I$ decays exponentially.
\end{proof}

\begin{intuition}
The number $R_0$ represents the average number of secondary infections
caused by a single infected individual during their entire infectious
period (of mean duration $1/\gamma$). The epidemic spreads only if each
patient infects more than one person on average ($R_0 > 1$).
\end{intuition}

\begin{definition}[SEIR Model]
\index{SEIR model}
The \textbf{SEIR} model adds an ``Exposed'' (incubation) compartment:
\[
  \frac{dS}{dt} = -\beta\,\frac{SI}{N}, \quad
  \frac{dE}{dt} = \beta\,\frac{SI}{N} - \sigma\,E, \quad
  \frac{dI}{dt} = \sigma\,E - \gamma\,I, \quad
  \frac{dR}{dt} = \gamma\,I,
\]
where $1/\sigma$ is the mean incubation period.
\end{definition}

\begin{keyformulas}
\textbf{Key Epidemiological Parameters}
\begin{align*}
  R_0 &= \frac{\beta}{\gamma} \quad \text{(SIR and SEIR)}, \qquad
  R_{\text{eff}}(t) = R_0 \cdot \frac{S(t)}{N} \\[4pt]
  \text{Final size:}\quad & S_\infty \text{ solves }
    S_\infty = S_0\, e^{-R_0(N - S_\infty)/N}
\end{align*}
\end{keyformulas}

\begin{example}
Calibrating the SIR model on influenza data from an English boarding
school (1978): 763 students, daily data over 14 days. Least-squares
fitting gives $\beta \approx 1.66$ and $\gamma \approx 0.44$,
yielding $R_0 \approx 3.78$.
\end{example}

\begin{minted}{python}
import numpy as np
from scipy.integrate import odeint
from scipy.optimize import minimize

def sir_model(y, t, beta, gamma, N):
    S, I, R = y
    dSdt = -beta * S * I / N
    dIdt = beta * S * I / N - gamma * I
    dRdt = gamma * I
    return [dSdt, dIdt, dRdt]

# Observed infected counts (daily)
data_I = np.array([1, 3, 8, 28, 76, 222, 293, 257, 237,
                   192, 126, 70, 28, 12])
t_data = np.arange(len(data_I))
N = 763

def objective(params):
    beta, gamma = params
    y0 = [N - 1, 1, 0]
    sol = odeint(sir_model, y0, t_data, args=(beta, gamma, N))
    return np.sum((sol[:, 1] - data_I)**2)

result = minimize(objective, [1.5, 0.5], method='Nelder-Mead')
\end{minted}

% ------------------------------------------------------------
\section{Population Dynamics: Lotka--Volterra}
\label{sec:lotka_volterra}
% ------------------------------------------------------------

\begin{definition}[Lotka--Volterra Model]
\index{Lotka--Volterra}
The predator--prey Lotka--Volterra model reads:
\[
  \frac{dx}{dt} = \alpha\,x - \beta\,x\,y, \qquad
  \frac{dy}{dt} = \delta\,x\,y - \gamma\,y,
\]
where $x(t)$ is the prey population, $y(t)$ the predator population,
and $\alpha, \beta, \gamma, \delta > 0$.
\end{definition}

\begin{proposition}[First Integral]
The Lotka--Volterra system admits the conserved quantity:
\[
  H(x,y) = \delta\,x - \gamma\,\ln x + \beta\,y - \alpha\,\ln y.
\]
Phase-plane trajectories are therefore closed curves around the
equilibrium $(\gamma/\delta,\, \alpha/\beta)$.
\end{proposition}

\begin{proof}
Compute $\frac{dH}{dt} = \delta\,\dot{x} - \frac{\gamma}{x}\,\dot{x}
+ \beta\,\dot{y} - \frac{\alpha}{y}\,\dot{y}$. Substituting the system
equations, all terms cancel: $dH/dt = 0$.
\end{proof}

\begin{attention}
The classical Lotka--Volterra model is structurally unstable: it has
no attracting limit cycle. Any perturbation changes the oscillation
amplitude permanently. For a more realistic model, add a carrying
capacity (logistic term) or a Holling-type functional response\index{Holling response}.
\end{attention}

\begin{example}
Historical data from the Hudson's Bay Company (hare and lynx pelts,
1845--1935) show phase-shifted oscillations characteristic of
predator--prey dynamics. Least-squares calibration yields periods
of approximately 10 years.
\end{example}

% ------------------------------------------------------------
\section{Traffic Flow Models}
\label{sec:traffic}
% ------------------------------------------------------------

\begin{definition}[LWR Model]
\index{LWR model}
The \textbf{Lighthill--Whitham--Richards} (LWR) model is a scalar
conservation law for the vehicle density $\rho(x,t)$:
\[
  \frac{\partial \rho}{\partial t} +
  \frac{\partial}{\partial x}\bigl[Q(\rho)\bigr] = 0,
\]
where $Q(\rho) = \rho\,v(\rho)$ is the flux and $v(\rho)$ is the
speed-density relation.
\end{definition}

\begin{proposition}[Greenshields Relation]
\index{Greenshields}
The Greenshields model assumes $v(\rho) = v_{\max}(1 - \rho/\rho_{\max})$,
giving a parabolic flux:
\[
  Q(\rho) = v_{\max}\,\rho\left(1 - \frac{\rho}{\rho_{\max}}\right).
\]
The capacity (maximum flux) is $Q_{\max} = v_{\max}\rho_{\max}/4$,
attained at $\rho = \rho_{\max}/2$.
\end{proposition}

\begin{remark}
Solutions of the LWR model can develop discontinuities (\emph{shocks}),
corresponding physically to traffic jam formation. The
Rankine--Hugoniot condition\index{Rankine--Hugoniot} gives the shock
speed: $s = [Q(\rho_R) - Q(\rho_L)]/(\rho_R - \rho_L)$.
\end{remark}

% ------------------------------------------------------------
\section{Introduction to Climate Modeling}
\label{sec:climate}
% ------------------------------------------------------------

\begin{definition}[Energy Balance Model]
\index{energy balance model}
The \textbf{Budyko--Sellers energy balance model}\index{Budyko} describes
the mean temperature $T$ of the Earth:
\[
  C\,\frac{dT}{dt} = (1 - \alpha(T))\,Q - \varepsilon\,\sigma\,T^4,
\]
where $C$ is the heat capacity, $Q$ the incoming solar flux,
$\alpha(T)$ the albedo (temperature-dependent via ice cover), and
$\varepsilon\,\sigma\,T^4$ the emitted infrared radiation.
\end{definition}

\begin{theorem}[Bifurcation and ``Snowball Earth'']
The Budyko--Sellers model generically admits three equilibria: a warm
state (little ice), a ``snowball'' state (fully glaciated Earth), and
an unstable intermediate equilibrium. When the solar flux $Q$ decreases,
a saddle-node bifurcation triggers an abrupt transition to the glaciated state.
\end{theorem}

\begin{intuition}
Ice albedo ($\sim 0.7$) is much higher than ocean albedo ($\sim 0.06$).
The colder it gets, the more ice forms, the more sunlight is reflected,
the colder it gets\ldots{} This is a positive feedback loop that can
lead to a runaway glaciation.
\end{intuition}

% ------------------------------------------------------------
\section{Model Comparison and Validation}
\label{sec:validation}
% ------------------------------------------------------------

\begin{definition}[Akaike Information Criterion (AIC)]
\index{AIC}
Let $\hat{L}$ be the maximum likelihood of a model with $k$ parameters.
The \textbf{Akaike Information Criterion} is:
\[
  \text{AIC} = -2\,\ln\hat{L} + 2k.
\]
The model with the smallest AIC is preferred.
\end{definition}

\begin{definition}[Bayesian Information Criterion (BIC)]
\index{BIC}
The \textbf{BIC} (or Schwarz criterion) is:
\[
  \text{BIC} = -2\,\ln\hat{L} + k\,\ln n,
\]
where $n$ is the number of observations. BIC penalizes complexity more
heavily than AIC when $n \geq 8$.
\end{definition}

\begin{proposition}[Cross-Validation]
\index{cross-validation}
$k$-fold \textbf{cross-validation} estimates the prediction error by
splitting the data into $k$ parts. The estimated error is:
\[
  \text{CV}(k) = \frac{1}{k}\sum_{j=1}^k \text{Error}_j,
\]
where $\text{Error}_j$ is the error of the model fitted on the other
$k-1$ parts and evaluated on part $j$.
\end{proposition}

\begin{attention}
A good fit to training data does not guarantee good predictive power.
\emph{Overfitting}\index{overfitting} occurs when the model captures noise
rather than signal. Always validate on independent data.
\end{attention}

\begin{keyformulas}
\textbf{Model Comparison Criteria}
\begin{align*}
  R^2 &= 1 - \frac{\sum_i (y_i - \hat{y}_i)^2}{\sum_i (y_i - \bar{y})^2}
    \quad \text{(coefficient of determination)} \\[4pt]
  \text{RMSE} &= \sqrt{\frac{1}{n}\sum_{i=1}^n (y_i - \hat{y}_i)^2}
    \quad \text{(root mean squared error)} \\[4pt]
  \text{AIC} &= -2\ln\hat{L} + 2k, \qquad
  \text{BIC} = -2\ln\hat{L} + k\ln n
\end{align*}
\end{keyformulas}

\begin{exercise}
Implement the SEIR model with COVID-19 parameters
($\sigma^{-1} \approx 5.2$ days, $\gamma^{-1} \approx 10$ days,
$R_0 \approx 2.5$). Simulate the dynamics for a city of $100{,}000$
inhabitants and study the effect of lockdown (reduced $R_0$).
\end{exercise}

\begin{exercise}
Calibrate a Lotka--Volterra model on the Hudson's Bay Company
hare--lynx dataset. Estimate parameters by nonlinear least squares
and assess fit quality via $R^2$ and RMSE.
\end{exercise}

\begin{exercise}
Compare SIR, SEIR, and SIR-with-vaccination models on an epidemic
dataset. Use AIC and BIC to determine the most parsimonious model.
\end{exercise}

\begin{exercise}
Solve the LWR model numerically with the Greenshields relation to
simulate traffic jam formation at a red light. Use a Godunov
scheme\index{Godunov scheme} and visualize the density
$\rho(x,t)$ in space and time.
\end{exercise}
